\begin{textsample}{POS Dim 3 – pa – Score 23.00 – pa/pa\_inf\_14.txt}  \label{ex:f3_pos_231}
3.1.8 Organização dos Espaços, Materiais e Tempos A trajetória histórica da Educação \textbf{Infantil} no Brasil foi sendo construída, entre outros aspectos, pelo assistencialismo, pelo descaso das políticas públicas, pela omissão de direitos, pela cisão entre o \textbf{cuidar} e o educar e pela separação entre ricos e pobres ( BRASIL, 2013a ).
As concepções que constituíram esses fatos e que estão atrelados a eles têm raízes profundas que atravessam o tempo e persistem até hoje ; atingem e colaboram na produção das concepções que povoam as mentalidades e os cenários educativos do século XXI.
E o que isso tem a ver com a organização dos espaços, materiais e tempos da Educação \textbf{Infantil}?
As velhas concepções atingem em \textbf{cheio} a temática que dá título a esse tema porque muitos ainda concebem o espaço como algo secundarizado, inerte, estéril ou que pouco contribuiu para a qualidade do atendimento.
Olhando com um pouco mais de cautela se pode verificar que o espaço isoladamente talvez não nos remeta a relações mais significativas com os processos educativos estabelecidos com as \textbf{crianças} e as \textbf{infâncias}, porém se o vincular ao tempo que é histórico e estabelece uma \textbf{íntima} relação com a produção cultural de materiais far -se-á uma ligação mais profícua entre a trilogia espaço-tempo-materiais.
Dessa forma se oferecer às \textbf{crianças} materiais atrativos para que elas criem e recriem possibilidades de uso, têm -se grandes chances de provocar significativas aprendizagens e com um pouco mais de diálogo, estudo e pesquisa pode se estar auxiliando as \textbf{crianças} a transformarem o movimento de suas vidas.
A isto se dá o nome de desenvolvimento.
A \textbf{criança} mobiliza -se em uma atividade, quando investe nela, quando faz uso de si mesma como de um recurso, quando é posta em movimento por móbeis que remetem a um \textbf{desejo}, um sentido, um valor.
A atividade possui, então, uma dinâmica interna.
Não se deve esquecer, entretanto, que essa dinâmica supõe uma troca com o mundo, onde a \textbf{criança} encontra metas desejáveis, meios de ação e outros recursos que não ela mesma ( CHARLOT, 2000, p. 55 ).
O autor chama atenção sobre a relação das \textbf{crianças} com o movimento.
Atitude própria e inerente ao ser \textbf{criança} e vivente no mundo.
O mundo é o espaço mediador, cenário provocador de novas e sempre desafiadoras possibilidades.
Nesse cenário o professor é o sujeito histórico cultural que, consciente dessas possibilidades e cronologicamente detentor de uma maior experiência cultural e histórica suscita possibilita o contato com novos espaços e materiais em situações que precisam ser estimuladoras.
Isso não quer dizer, no entanto, que o professor está sempre a controlar tudo, como se assumisse ser o senhor do tempo para ditar quando e como se deve iniciar e parar as coisas.
O professor sutilmente deve \textbf{atentar} -se e buscar intimidade com as \textbf{crianças} e com as relações de significância que elas fazem com os objetos, pessoas e situações.
Ele deve ver, \textbf{ouvir}, perceber, o que as \textbf{crianças} dizem como se deslocam no espaço, como dão vida aos materiais e o modo como os ressignificam, como estabelecem relações com pessoas situações, objetos, elementos da natureza.
Sobre o que \textbf{conversam} o que imaginam e sentem.
Dessa forma, o espaço deve oferecer materiais acessíveis e \textbf{interessantes} para o olhar das \textbf{crianças}, tornando -se local aprazível e acolhedor da intimidade do ser \textbf{criança}, sendo que intimidade não quer dizer que lá só haja lugar para o que é “ conhecido ”.
É necessário que as \textbf{crianças} estejam em contato com novos e inusitados materiais que lhes proporcionem exploração e surpresas, como caixas, carreteis, pedaços de madeiras, tronco de árvores, folhas secas, tendas, são inúmeras as possibilidades de materiais que podem colaborar para que o espaço vá se transformando em lugar.
O espaço se projeta ou se imagina ; o lugar se constrói.
Constrói -se a partir do fluir da vida, das relações que ali são travadas e a partir do espaço como suporte ; o espaço, portanto, está sempre disponível e disposto para converter -se em lugar, para ser construído ( AGOSTINHO, 2003, p. 1 ).
Essa forma de ver o espaço, revelada pela autora, \textbf{convida} -nos a rever nossas posturas e o modo como as práticas são conduzidas e reproduzidas no espaço escolar ; as próprias DCNEI ( BRASIL, 2010a ), em seu Art. 8º § 1º, chama atenção para alguns aspectos relacionados aos objetivos da proposta pedagógica no que tange à organização de materiais, espaços e tempos e chama atenção dentre outros fatores em assegurar a integralidade, e indissociabilidade entre o \textbf{cuidar} e o educar, o reconhecimento das especificidades que orientam o desenvolvimento das \textbf{crianças}, a mobilidade nos espaços, além da acessibilidade de espaços, materiais objetos, \textbf{brinquedos}, para \textbf{crianças} com deficiência, transtornos globais de desenvolvimento e altas habilidades/superdotação e apropriação, reconhecimento e valorização das múltiplas culturas que contribuíram para a formação do povo brasileiro ( BRASIL, 2010a ).
Como se pode perceber, os espaços, tempos e materiais devem ser pensados no contexto da Educação \textbf{Infantil} de forma \textbf{lúdica} dando oportunidade para que as \textbf{crianças} possam exercitar sua criatividade, exploração e descobertas.
O espaço ele não é só um pano de fundo para as aprendizagens.
O espaço interfere diretamente na aprendizagem das \textbf{crianças}, ele é inclusive entendido, em algumas experiências educacionais, como interlocutor, como educador, inclusive, porque ele desafia, porque ele instiga as \textbf{crianças} à exploração, ao movimento, a produção de linguagens ( GUIMARÃES, 2010, n.p. ).
Perceber o espaço como algo vivo e que dialoga com o processo educacional é entendê -lo como interlocutor ; da mesma forma, pensar em estratégias que possam fazer com que a \textbf{curiosidade} e a criatividade das \textbf{crianças} estejam sempre sendo estimuladas por esse ambiente é criar possibilidades novas e imprevisíveis para que elas se desenvolvam com qualidade e prazer.
Na maioria das vezes o que se observa em nossas práticas com a Educação \textbf{Infantil} é que o professor ainda está preso às velhas práticas promovendo, quase sempre, atividades \textbf{repetidas} e que priorizam grandes grupos de mesma faixa \textbf{etária} ; todos fazem ao mesmo tempo as mesmas coisas limitando \textbf{rotinas} e metodologias, fixando procedimentos e cerceando possibilidades de \textbf{crescimento} humano-intelectual o que é mais comprometedor ainda.
Alargar o olhar para repensar a organização do espaço requer uma mudança ampla nos caminhos da escola, mudança esta que deve partir de um movimento coletivo da escola e não apenas de mentes que pensam de forma isolada.
O espaço deve proporcionar além do acesso e disponibilidade de objetos de diferentes formas, \textbf{texturas} e tamanhos, promover a interação entre as \textbf{crianças}, não só de mesma faixa \textbf{etária}, mas também de faixas \textbf{etárias} diferentes de modo que semanalmente as \textbf{crianças} possam interagir em grupos diferentes dos que aqueles estabelecidos em sua turma de origem.
Vale ressaltar que todas as reflexões aqui feitas devem estar em consonância com a proposta pedagógica curricular da \textbf{instituição} enquanto projeto pedagógico maior que rege e dá direção às intervenções estabelecidas sem deixar de problematizar em sua constituição a complexa relação histórica, política, social, econômica e cultural das relações humanas imbricadas.
O modo como as concepções de \textbf{criança} e \textbf{infância} dialogam promoverá ou não a interação entre crianças- ambientes e aprendizagens ( \textbf{OLIVEIRA}, 2012 ).
Para a autora, a construção de ambiente de convivência e aprendizagem na Educação \textbf{Infantil} necessita ser analisada sob a perspectiva de diferentes dimensões que precisam dialogar entre si para ultrapassar uma lógica limitada que nos conduz ao olhar ingênuo de que o espaço está restrito apenas ao aspecto espacial, mas se amplia em outras direções envolvendo : a ) Dimensão interacional – diferentes perspectivas de interação entre \textbf{crianças} de diferentes faixas \textbf{etárias} e destas com \textbf{adultos} ; b ) Dimensão física – considera o espaço como elemento educador.
Assume, portanto perspectiva de assegurar, estimular, renovar e planejar sempre novas formas de estabelecer formas diversas de organização ; por isso ele precisa ser pensado inclusive pelas próprias \textbf{crianças}.
Exigirá do professor também um olhar sensível para perceber que ali naquele espaço elas se sentem melhor do que o espaço originalmente pensado pelo professor, por exemplo ; c ) Dimensão temporal – implica considerar a qualidade do tempo vivido pelas \textbf{crianças} na \textbf{instituição}.
O que Arroyo ( 2011 ) denomina “ direitos a tempos-espaços de um justo e digno viver ” ; implica tornar o tempo produtivo, com etapas distintas e sempre renovadas ; isso também implica protagonizar o olhar das \textbf{crianças} em relação à ordem estabelecida para as atividades e ao tempo de duração de cada uma delas. d ) Dimensão funcional – considera a real função e significado dos espaços e suas possibilidades.
Estão de fato a serviço das \textbf{crianças}?
Contemplam na prática suas necessidades de segurança, mobilidade, descanso, alimentação, criação e exploração?
Podem contemplar mais de uma função de acordo com a motivação e ideias socializadas por professores e \textbf{crianças}?
Estas são algumas perguntas que podemos nos fazer como profissionais, mas podem ser feitas às \textbf{crianças} de modo a saber como estão se sentindo naquele ambiente e quais suas impressões sobre eles.
Essas dimensões ajudam a pensar um pouco mais sobre o espaço como \textbf{partícipe} do processo de aprendizagem no sentido de ser entendido como comunicador e termômetro para os índices de qualidade no ambiente escolar.
Sempre lembrando que, para além do espaço físico e das estruturas aparentemente rígidas em que convencionalmente se entende o tempo e os materiais, existe a humanidade que reside em cada \textbf{criança}, as interações subjetivas que superam e burlam a lógica que se pensa deter e a complexidade da existência humana que busca incessantemente mudanças e transformações.
Se a educação não se der nesse movimento, não é de fato educação.

% matched lemmas: adulto, atentar, brinquedo, cheio, conversar, convidar, crescimento, criança, cuidar, curiosidade, desejo, etário, infantil, infância, instituição, interessante, lúdico, oliveira, ouvir, partícipe, repetir, rotina, textura, íntimo
\end{textsample}
